% !TEX root = ../matan.tex

\section{Теорема сравнения, признак Даламбера}

\subsection*{Признак сравнения}

\begin{Theorem}{Признак сравнения}{}
	Пусть $\sum_{k=1}^{\infty} b_k$ --- ряд с неотрицательными членами ($b_k \ge 0$),
	а $\sum_{k=1}^{\infty} a_k$ --- произвольный ряд. Тогда:
	\begin{enumerate}
		\item если $|a_k| \le b_k$ для всех $k$ и ряд $\sum_{k=1}^{\infty} b_k$ сходится,
		то ряд $\sum_{k=1}^{\infty} a_k$ сходится;
		\item если $a_k \ge b_k$ для всех $k$ и ряд $\sum_{k=1}^{\infty} b_k$ расходится,
		то ряд $\sum_{k=1}^{\infty} a_k$ расходится.
	\end{enumerate}
\end{Theorem}

\begin{proof}
	\emph{(1)} Так как $\sum b_k$ сходится, по критерию Коши для ряда:
	для любого $\varepsilon > 0$ найдётся $N(\varepsilon)$ такое, что для всех
	$n \ge N(\varepsilon)$ и $m \in \mathbb{N}$
	\[
	\left|\sum_{k=n+1}^{n+m} b_k\right| < \varepsilon.
	\]
	Поскольку $b_k \ge 0$, левая часть равна $\sum_{k=n+1}^{n+m} b_k$.
	Из условия $|a_k| \le b_k$ и неравенства треугольника:
	\[
	\left|\sum_{k=n+1}^{n+m} a_k\right|
	\;\le\; \sum_{k=n+1}^{n+m} |a_k|
	\;\le\; \sum_{k=n+1}^{n+m} b_k
	\;<\; \varepsilon.
	\]
	По критерию Коши ряд $\sum a_k$ сходится. В частности, из этих же неравенств
	видно, что сходится и $\sum |a_k|$, то есть $\sum a_k$ сходится абсолютно.
	
	\emph{(2)} Так как $b_k \ge 0$ и $\sum b_k$ расходится, его частичные суммы
	стремятся к $+\infty$:
	\[
	\sum_{k=1}^{n} b_k \;\xrightarrow{n \to +\infty}\; +\infty.
	\]
	Из условия $a_k \ge b_k \ge 0$ получаем
	$\sum_{k=1}^{n} a_k \ge \sum_{k=1}^{n} b_k \to +\infty$,
	поэтому ряд $\sum a_k$ расходится.
\end{proof}

\subsection*{Предельный признак сравнения}

\begin{Theorem}{Предельный признак сравнения}{}
	Пусть $a_k \ge 0$, $b_k > 0$ и существует конечный положительный предел
	\[
	\lim_{k\to\infty} \frac{a_k}{b_k} = c, \qquad 0 < c < +\infty.
	\]
	Тогда ряды $\sum a_k$ и $\sum b_k$ сходятся или расходятся одновременно.
\end{Theorem}

\begin{proof}
	Так как $\frac{a_k}{b_k} \to c$ и $c > 0$, начиная с некоторого $k_0$
	выполнено $\frac{c}{2} \le \frac{a_k}{b_k} \le \frac{3c}{2}$, то есть
	\[
	\frac{c}{2}\,b_k \;\le\; a_k \;\le\; \frac{3c}{2}\,b_k \qquad (k \ge k_0).
	\]
	Из правого неравенства по признаку сравнения сходимость $\sum b_k$ влечёт
	сходимость $\sum a_k$; из левого --- сходимость $\sum a_k$ влечёт
	сходимость $\sum b_k$.
\end{proof}

\subsection*{Признак Даламбера}

\begin{Theorem}{Признак Даламбера}{}
	Пусть $\sum_{k=1}^{\infty} a_k$ --- произвольный ряд, $|a_k| > 0$ для
	всех $k$. Тогда:
	\begin{enumerate}
		\item если существуют $q < 1$ и $k_0 \in \mathbb{N}$ такие, что
		\[
		\left|\frac{a_{k+1}}{a_k}\right| \le q
		\quad \text{для всех } k \ge k_0,
		\]
		то ряд $\sum_{k=1}^{\infty} a_k$ сходится;
		\item если существуют $q > 1$ и $k_0 \in \mathbb{N}$ такие, что
		\[
		\left|\frac{a_{k+1}}{a_k}\right| \ge q
		\quad \text{для всех } k \ge k_0,
		\]
		то ряд $\sum_{k=1}^{\infty} a_k$ расходится.
	\end{enumerate}
\end{Theorem}

\begin{proof}
	Идея: сравнение с геометрической прогрессией.
	
	\emph{(1)} Запишем частичную сумму:
	\[
	S_n = \sum_{k=1}^{n} a_k
	= \underbrace{\sum_{k=1}^{k_0-1} a_k}_{\text{фикс. конечная сумма}}
	+ \sum_{k=k_0}^{n} a_k.
	\]
	Первая часть --- фиксированное конечное число, на сходимость не влияет.
	Исследуем хвост $\sum_{k=k_0}^{n} a_k$.
	
	Из условия $|a_{k+1}| \le q \cdot |a_k|$ при $k \ge k_0$ по индукции:
	\[
	|a_{k_0 + m}| \le |a_{k_0}| \cdot q^m, \qquad m \ge 0.
	\]
	Следовательно,
	\[
	\sum_{k=k_0}^{n} |a_k|
	= \sum_{m=0}^{n-k_0} |a_{k_0+m}|
	\le |a_{k_0}| \sum_{m=0}^{n-k_0} q^m
	\le \frac{|a_{k_0}|}{1-q} < +\infty,
	\]
	так как $q < 1$. Частичные суммы $\sum_{k=k_0}^{n} |a_k|$ монотонно возрастают
	и ограничены, значит ряд $\sum_{k=k_0}^{\infty} |a_k|$ сходится. По признаку
	сравнения ряд $\sum_{k=k_0}^{\infty} a_k$ сходится абсолютно, а значит,
	сходится. Вместе с конечной начальной частью ряд $\sum_{k=1}^{\infty} a_k$ сходится.
	
	\emph{(2)} Из условия $|a_{k+1}| \ge q \cdot |a_k|$ при $k \ge k_0$
	по индукции:
	\[
	|a_{k_0 + m}| \ge |a_{k_0}| \cdot q^m \;\xrightarrow{m \to \infty}\; +\infty,
	\]
	так как $q > 1$. Значит, $a_k \not\to 0$. По необходимому условию сходимости
	ряд $\sum a_k$ расходится.
\end{proof}

\begin{Remark}{Случай $q = 1$ неопределён}{}
	Если отношение $\left|\frac{a_{k+1}}{a_k}\right|$ стремится к $1$, то ни
	условие (1), ни условие (2) не выполнено ни для какого $q < 1$ или $q > 1$:
	признак Даламбера ответа не даёт.
	
	Пример: для $a_k = \dfrac{1}{k}$
	\[
	\frac{a_{k+1}}{a_k} = \frac{k}{k+1} \;\xrightarrow{k\to\infty}\; 1,
	\]
	и ряд $\sum \dfrac{1}{k}$ \emph{расходится}. Для $a_k = \dfrac{1}{k^2}$
	\[
	\frac{a_{k+1}}{a_k} = \frac{k^2}{(k+1)^2} \;\xrightarrow{k\to\infty}\; 1,
	\]
	но ряд $\sum \dfrac{1}{k^2}$ \emph{сходится}. Один и тот же предельный
	результат соответствует и сходящемуся, и расходящемуся ряду.
\end{Remark}